% Generated by records_reorg.py assemble (records-reorg/); edit the data and re-run, not this file.
\subsection{The unfolded cross section and the comparison with MINERvA}\label{recsec:twod-result}

% Part 3 (DECISION-20261119c; REWRITE-RULES-20261119.md rule 5), director-helper-3: detail the condensed Part 3 no longer prints, moved verbatim from sec_results.tex @ draft 47c318849, one move-manifest row per passage (H3-P3-01..H3-P3-09 in history-cuts/move-manifest-h3-20261117.tsv). None of these passages carries a citation. [records reorg] former heading: \section{Reproduction detail moved from the note (Part 3)} [records reorg, not printed] This section holds supporting detail on the two-dimensional result that the compressed Part~3 no longer prints, as the note printed it.

\paragraph{The shapes of the projections.} The $\pT$ projection vanishes at $\pT=0$ and peaks broadly at $\pT\approx0.55$--$\SI{0.70}{GeV/c}$, reflecting the $Q^{2}\lesssim\SI{1}{GeV^2}$ scale of few-GeV scattering through $\pT^{2}\simeq Q^{2}\,E_\mu/E_\nu$, with $E_\mu/E_\nu\sim0.6$ for these forward muons. It falls by more than two orders of magnitude toward the \SI{4.5}{GeV/c} edge. The $\pPar$ projection rises from the \SI{1.5}{GeV/c} threshold to a broad $\pPar\approx3$--$\SI{4.5}{GeV/c}$ maximum (the NuMI medium-energy flux peak at $E_\nu\approx\SI{6}{GeV}$), then falls by two orders of magnitude by $\pPar\approx\SI{20}{GeV/c}$.
% [records reorg] former heading: \subsection{The shapes of the projections (from Part 3, the double-differential cross section)}
% From sec_results.tex:9 @ draft 47c318849, between "The shapes are those expected ..." and "The $\theta_\mu<20^{\circ}$ acceptance ..."

\paragraph{The error bars of the projections.} The bars on this result in Fig.~\ref{fig:xsec} are the square root of summed squared event weights, combined in quadrature with the propagated completeness-ratio bin error (ratio one by construction).
% [records reorg] former heading: \subsection{The error bars of the projections (from the caption of Fig.~\ref{fig:xsec})}
% From sec_results.tex:16 @ draft 47c318849, the clause after "Bars on this result: weight statistics only, not the covariance;" (verbatim clause)

\paragraph{Standardized differences and the slice comparisons.} Per-bin differences standardized by the published uncertainty (Fig.~\ref{fig:pull}) are centered (mean $\pullMean$), with standard deviation $\pullRMSr$, well inside the published uncertainties. Figures~\ref{fig:ptslices}--\ref{fig:pzslices} compare the unfolded result with \emph{simulated truth}, the MC prior, in the published panel format, with the published points for comparison.
% [records reorg] former heading: \subsection{Standardized differences and the slice comparisons (from Part 3, the comparison to the published measurement)}
% From sec_results.tex:83 @ draft 47c318849 (whole paragraph); Table~\ref{tab:compare} prints the mean and standard deviation

\paragraph{Eigenmode detail of the published-covariance $\chi^{2}$.} The largest mode has a paper-covariance-standardized component of magnitude $4.4$, and the next nine $3.1$--$4.0$. The 10 smallest-$\lambda$ modes carry \SI{8.7}{\percent} of the total $\chi^{2}$. Of the five largest contributors, one lies in the smallest-$\lambda$ tenth of the spectrum. Of the six leading modes, the five in the highest-$\pPar$ bins are all poorly determined directions ($\lambda/\lambda_{\max}\lesssim10^{-8}$). The sixth is a moderate-$\lambda$ mode ($\lambda/\lambda_{\max}\approx6\times10^{-5}$) and lies below the highest $\pPar$ bins (Fig.~\ref{fig:tensionmodes}).
% [records reorg] former heading: \subsection{Eigenmode detail of the published-covariance $\chi^{2}$ (from Part 3, where the tension lies)}
% From sec_results.tex:130 @ draft 47c318849, after "The eigenmode decomposition (Fig.~\ref{fig:tension}) spreads it over many modes", and the sentences before "No single ... band"
% D2 (crossdoc audit 2026-11-15): per-period-benchmark-20261106/receipts/read/readout.json:61 abs_z 4.746664876937429 (the 4.8 rounded the rounded 4.75, :62); the Supplement prints pp-mode-max, 4.7
% 2026-11-25 (pT-edge option (i); Joseph's wording, incl. "The largest mode has" for the reorg's fragment): on the published
% pT edges, regenA A5 (ibu-chi2-20261123/regenA5/out/readout_e.json, #1330): largest_abs_z 4.394, next nine 3.103--4.031,
% share10 8.71%, top5_in_smallest_tenth 1, six_split 5/1. A5's OLD reproduced 4.7, 3.4--3.8, 9.6% and two. Bound by
% check_sec4_receipts.py (rows "A5 ...").

\paragraph{The truth denominator of the slice figures.} The analysis-MC truth denominator uses the unfolded result's flux, nucleon-count and efficiency normalization.
% [records reorg] former heading: \subsection{The truth denominator of the slice figures (from the caption of Fig.~\ref{fig:ptslices})}
% From sec_results.tex:113 @ draft 47c318849

\paragraph{From the two-dimensional cross section (\S\ref{sec:results}).} The comparison leaves out the generator-shape component $V_{\mathrm{prior}}$ (\S\ref{sec:uq-2d-scope}), because the published uncertainty has no such component either. With $V_{\mathrm{prior}}$, the per-bin median is \SI{\qrCmodelMedian}{\percent}. The systematic-universe block accounts for a median 0.87 of the per-bin variance and 0.83 of the integral's variance, and the run-period block (\S\ref{sec:uq-combine}) raises the per-bin median from \SI{\uqSFiveMedian}{\percent} to \SI{\uqMedian}{\percent}. Our flux band carries the per-universe $\nu$--$e$ and inverse-muon-decay constraint weights of the MINERvA Analysis Toolkit. The statistical and ML blocks together contribute under a tenth of the combined $\sqrt{\mathrm{Tr}\,C}$. Every block is computed with the LightGBM estimator of the quoted central value (\S\ref{sec:method-config}). Relative uncertainties are taken with respect to that central value, which is the universe sweep's central value at training seed~42. With the covariance of this result alone, $\chi^{2}$ is computed only in the covariance's well-constrained subspace, with a truncated spectral pseudo-inverse. The per-bin pulls and the truncated result are therefore quoted together, and an unconstrained pseudo-inverse is not used as a goodness-of-fit test on its own. The regularization scan (App.~\ref{sec:rank}) was made for the exact-split cross-check on an earlier covariance and was not repeated for the quoted one. Restricted to the 73 best-measured directions, $\chi^{2}=136$ ($\chi^{2}/205=0.66$; the divisor is the number of reported bins, not the number of directions kept). No single $\pT$ or $\pPar$ band carries more than \SI{15.2}{\percent} of the total. A value near one does not by itself indicate reproduction. What the benchmarks leave open is stated in \S\ref{rlim:2d}. \emph{After, in \S\ref{sec:results-benchmarks}, the note's ``A benchmark specified before it ran unfolds three groups of run periods separately and compares combinations of them with the published result in the directions of the published covariance.'' ``\lrval{pp-combined} In the other \lrval{pp-nkbar} poorly determined directions of the published covariance, the combination of the other two run-period groups is farther from the published result than the three-group combination.''} It is farther by \lrval{pp-kbar-delta} per direction, a fraction \lrval{pp-kbar-frac} of the three-group distance, beyond noise estimated from \lrval{pp-nrep} replicas. The two-dimensional estimator with deterministic training, the five-observable estimator's training mode, agrees with the two-dimensional central value within its paired noise, narrowly: $T=\fvval{c1-T}$ over \fvval{c1-nrep} paired replicas, \SI{\fvval{c1-below-pct}}{\percent} below its 95th-percentile threshold of \fvval{c1-q95}. Our result tracks the data far more closely than the model does (the earlier note's figure of the model-ratio maps). The published analysis~\cite{MINERvA:2021owq} also compares with NuWro~\cite{Golan:2012wx}, GiBUU~\cite{Buss:2011mx} and NEUT~\cite{Hayato:2009zz}. OmniFold's per-event reweighted truth ensemble can be compared with any such prediction in \emph{any} projection of the classifier features without re-unfolding. An observable outside that space requires a new unfolding, as in the 3D extension (\S\ref{sec:3d}), which adds $\Eavail$ as a classifier feature.
% [records reorg] former heading: \subsection{From the two-dimensional cross section (\S\ref{sec:results})}
% From sec_results.tex, paragraph 3
% [lead-in, Option B] \emph{After, in \S\ref{sec:results-uq}, the note's ``Table~\ref{tab:uq} gives the budget over the 205 reported bins, constructed as described in \S\ref{sec:systematics}.'' ``Our own combined budget, without the published covariance, differs from the published total in both directions.'' ``Its per-bin median, \SI{\uqMedian}{\percent}, is smaller than the published \SI{\uqPaper}{\percent}.'' ``Its value on the integrated cross section, \SI{\uqIntegral}{\percent}, is larger than the published \SI{\uqPaperIntegral}{\percent} (the width-weighted total $\sigma$ of App.~\ref{sec:chisq}).''}
% From sec_results.tex, paragraph 4
% [lead-in, Option B] \emph{After, in \S\ref{sec:results-uq}, the note's ``The largest band is the flux, as in the publication (the earlier note's figure of the uncertainty bands), and it accounts for part of the difference between the per-bin medians: its median is \SI{\fluxBand}{\percent}, against \SI{4.01}{\percent} for the published flux block.'' ``The public release cannot establish the origin of this difference, for which the FHC-only constraint file remains the leading, unconfirmed hypothesis of this analysis.''}
% From sec_results.tex, paragraph 7
% [lead-in, Option B] \emph{After, in \S\ref{sec:results-compare}, the note's ``The latter, $\chi^{2}/\mathrm{ndf}=\chiCombined$, double-counts the flux, GENIE and MINOS systematics that the two results share.''}
% From sec_results.tex, paragraph 8
% [lead-in, Option B] \emph{After, in \S\ref{sec:results-tension}, the note's ``With the published uncertainties as a descriptive scale, the central values have a pull standard deviation of $\pullRMS$ (a diagonal-only $\chi^{2}/\mathrm{ndf}$ would be $0.29$), whereas $\chi^{2}$ per bin with the full published covariance is $\chiPaper$.'' ``The eightfold increase lies entirely in the off-diagonal structure of the published covariance.'' ``It is a difference in \emph{shape}.'' ``Normalization does not explain it ($\chi^{2}_{\mathrm{norm}}=0.04$ against $\chi^{2}_{\mathrm{shape}}=465$).'' ``The eigenmode decomposition (Fig.~\ref{fig:tension}) spreads it over many modes.''}
% From sec_results.tex, paragraph 12
% [lead-in, Option B] \emph{After, in \S\ref{sec:results-benchmarks}, the note's ``Applied to this analysis's inputs, ten-iteration binned IBU and the two-dimensional OmniFold result give descriptive $\chi^{2}/\mathrm{ndf}$ values of \fvval{c2c1-dibu} and \fvval{c2c1-der}, respectively, relative to the published result under its covariance.'' ``IBU gives \fvval{c2c1-dibu-mid-lo}--\fvval{c2c1-dibu-mid-hi} for two to fifteen iterations, \fvval{c2c1-dibu-k20} at twenty and \fvval{c2c1-dibu-k1} at one.'' ``At fixed IBU settings, we replaced the classifier-refined reconstructed target with the binwise data-minus-background histogram.'' ``Evaluated with the published covariance, the descriptive $\chi^{2}/\mathrm{ndf}$ between the recorded OmniFold and IBU results changed from \fvval{ibusd-k5-ref} to \fvval{ibusd-k5-sig} at five iterations and from \fvval{ibusd-k10-ref} to \fvval{ibusd-k10-sig} at ten iterations.'' ``These comparisons use the ill-conditioned published covariance and lack the cross-covariance required for a statistical comparison of results that share data.''}
% From sec_results.tex, paragraph 13
% From sec_results.tex, paragraph 14
% From sec_results.tex, paragraph 15
% [lead-in, Option B] \emph{After, in \S\ref{sec:results-model}, the note's ``GENIE 2.12.6~\cite{Andreopoulos:2009rq} with MINERvA Tune v1 under-predicts the published total cross section over the 205 reported bins by $\sim\SI{9}{\percent}$ ($\sigma_{\mathrm{tune}}/\sigma_{\mathrm{data}}=0.91$), and the deficit is concentrated in the low-$\pT$ peak.'' ``Figure~\ref{fig:xsec} compares the prediction from the published ancillary files (dashed) with both the published data and this result.'' ``Under this covariance the quadratic distance to that prediction is $\chi^{2}/\mathrm{ndf}=33.0$ for the published data and $31.3$ for this result ($26.5$ for the exact-split cross-check).'' ``These are descriptive distances.'' ``This covariance gives no calibrated significance for a generator comparison (\S\ref{sec:uq-2d-scope}).''}


% H3-P3-09: sec_tension.tex (the former App. sec:tension), moved whole and byte-identical from the note's Part 3 (part3_reproduction.tex @ draft 47c318849). Its numbers are for the exact-split cross-check, the earlier central value; the note keeps the quoted value's decomposition (sec:results-tension).
\subsection{Anatomy of the paper-covariance tension}
\label{sec:tension}

This derivation supports \S\ref{sec:results-tension}. Its numbers are for the exact-split cross-check, the quoted 2D central value when the derivation was made (\S\ref{sec:method-config}). For the quoted LightGBM central value, \S\ref{sec:results-tension} gives the same decomposition from the same code: a smaller distance ($\chiPaper$), an $8\times$ rather than $10\times$ off-diagonal inflation, \SI{9.6}{\percent} rather than \SI{3}{\percent} in the ten smallest-$\lambda$ modes, and its leading poorly determined modes in the highest-$\pPar$ bins, while its moderate-$\lambda$ modes lie below the highest-$\pPar$ bins.

\paragraph{The paradox.} The central values agree well. For the exact-split cross-check, $\sigma_{\tot}^{\ours}/\sigma_{\tot}^{\paper} = 1.011$, the median bin ratio is $1.006$ and the per-bin pull standard deviation against the paper covariance is only $0.598$; for the quoted value they are $\ratioTot$, $\medianBinRatio$ and $\pullRMS$ (\S\ref{sec:pull}). A diagonal-only $\chi^{2}$, ignoring all off-diagonal correlations, is therefore % N4 (the director, 2026-11-16): the three exact-split values are values.tex's comments' (integral ratio 1.0113, median 1.0064, pull RMS 0.598); the quoted value's are \ratioTot, \medianBinRatio and \pullRMS (values.tex:10, :14, :19).
\begin{equation}
\chi^{2}_{\diag}/\mathrm{ndf}
= \frac{1}{n}\sum_{i}\frac{\Delta_{i}^{2}}{C_{ii}}
= \mathrm{SD}(\mathrm{pull})^{2} + \mathrm{mean}(\mathrm{pull})^{2}
\approx 0.366,
\end{equation}
versus full-covariance $\chi^{2}/\mathrm{ndf}=3.661$. The $10\times$ inflation comes entirely from published \texttt{TotalCov}'s off-diagonal and small-eigenvalue structure.
% (\texttt{diagnose\_tension.py})

\paragraph{Eigenmode decomposition.} Diagonalizing the reported-block paper covariance, $\mathbf{C}=\sum_{k}\lambda_{k}\mathbf{u}_{k}\mathbf{u}_{k}\T$ ($\lambda_{1}\!\ge\!\cdots\!\ge\!\lambda_{n}$), and projecting $c_{k}=\mathbf{u}_{k}\T\boldsymbol{\Delta}$,
\begin{equation}
\chi^{2}=\boldsymbol{\Delta}\T\mathbf{C}^{+}\boldsymbol{\Delta}
=\sum_{k}\frac{c_{k}^{2}}{\lambda_{k}},
\qquad
\text{``eigen-pull''}_{k}\equiv\frac{c_{k}}{\sqrt{\lambda_{k}}},
\end{equation}
reproduces the pseudo-inverse $\chi^{2}=750.49$ to $3\times10^{-9}$. The tension is neither a few pathological null directions nor the absolute smallest eigenvalues. The largest contributions lie at $\lambda_{k}/\lambda_{\max}\sim10^{-4}$--$10^{-5}$, in the lower-middle spectrum, $\sim7$--$8$ decades above the condition floor ($\lambda/\lambda_{\max}\sim10^{-12}$). Their standardized-component magnitudes $5$--$6$ are descriptive coordinates in the published basis, without calibrated significance for independent results. The ten \emph{smallest}-$\lambda$ modes carry only $3\%$ of the $\chi^{2}$ and about $90$ modes are needed to reach $90\%$: the $\chi^{2}$ is broadly distributed, not localized to the ill-conditioned tail. A scan of the pseudo-inverse cutoff excludes a numerical artifact: the $73$ best-determined directions ($\lambda_{k}/\lambda_{\max}>10^{-4}$) give $\chi^{2}=291$ ($\chi^{2}/205=1.42$), and $139$ directions ($>10^{-6}$) give $\chi^{2}=572$ ($\chi^{2}/205=2.79$), each divided by the number of reported bins, not by the number of directions kept. The rank-truncation scan of \S\ref{sec:chisq} rises smoothly with no late jump, so the tension does not depend on the last few least-determined modes.
% Revision D' 2026-11-06 (referee reads of D, both ACCEPT; Sonnet and the interpreter: print the raw chi2 and the directions kept): docs/orchestration/receipts/2d-lgbm42-20261011/tension_report_exact_control.txt:55-56 print "1e-06  2.791  139" and "1e-04  1.419  73". Arithmetic: 2.791 x 205 = 572.2, in [572.05, 572.26], 572 at three figures; 1.419 x 205 = 290.9, in [290.8, 291.0], 291.
% Director-added wording correction 2026-11-06, for Joseph's review before merge (not among the two he approved): the truncated chi2 was printed as chi2/ndf, but its divisor is the 205 reported bins, not the directions kept. 2d-unfolding/diagnose_tension.py @ c3dbe30f4 (main): :110 n = int(mask.sum()), the 205 reported bins; :183 and :193 divide every truncated chi2 by n (the same divisor at 56e30b9c9 :108/:181/:191 and at 735c909da). Receipt docs/orchestration/receipts/2d-lgbm42-20261011/tension_report_exact_control.txt:55-56 ("1e-06  2.791  139", "1e-04  1.419  73"). No number changed.

\paragraph{Shape contribution.} Splitting $\boldsymbol{\Delta}$ along $\mathbf{x}^{\paper}$ and its orthogonal remainder gives $\chi^{2}_{\mathrm{norm}}=0.10$ for scale, $\chi^{2}_{\text{shape}}=750.1$ for shape, and cross term $0.28$; the split is not $\mathbf{C}^{+}$-orthogonal. The $+1.1\%$ normalization offset is statistically irrelevant to the tension; it is a correlated shape difference, consistent with Jacobian shape-only $\chi^{2}/\mathrm{ndf}=3.60$ (\S\ref{sec:shape}).

\paragraph{Location in the $(\pT,\pPar)$ plane.} Mapping dominant eigenvectors and per-bin contributions $\Delta_{i}(\mathbf{C}^{+}\boldsymbol{\Delta})_{i}$ to the $14\times16$ grid (numbers for the exact-split cross-check; Figs.~\ref{fig:tension} and~\ref{fig:tensionmodes} show the LightGBM central value) locates tension on the low-$\pT$ peak ridge ($\pT\lesssim0.4$ GeV/c, $\pT$-bins 2/7/10 carry $16/11/12\%$) and low-$\pPar$ near-edge cells ($\pPar$-bins 0/1/8). Carrier modes oscillate in sign across the peak. They form neither vertical $\pT$-columns characteristic of a flux $1/\Phi(\pT)$ normalization error nor the $\theta_{\mu}<20^{\circ}$ acceptance boundary. Flux and Muon\_Energy, the largest bands and only named cross-band correlation (Bashyal, \S\ref{sec:rank}), dominate there. Tight bin-to-bin correlations leave little freedom to absorb a coherent $\sim1$--$2\%$ shape difference, so a small OmniFold-vs-IBU shape difference becomes a large $\chi^{2}$.\footnote{The pipeline uses one global $\mathrm{POT}_{\mathrm{data}}/\mathrm{POT}_{\mathrm{MC}}$ ratio; per-playlist ratios span $0.1707$--$0.2371$. Sensitivity at $1$--$2\%$ motivated a direct mixture check: $R_{\mathrm{glob}}\sum_p \mathrm{MC}_p$ versus $\sum_p R_p\,\mathrm{MC}_p$ over twelve playlists changes reco shape by at most $0.07\%$ in $\pT$, $0.14\%$ in $\pPar$, and $0.03\%$ on the low-$\pT$ ridge. The residual $+0.12\%$ is normalization. Nearly shape-identical playlists make even large mixture reweighting almost pure normalization, leaving the statement above unaffected at the $14$--$30\times$ level.}
%

\paragraph{A methodological component (estimator and iteration).} Changing the GBDT backend on the same data and normalization ($\Sigma\sigma$ within $0.1\%$) shifts the $\chi^{2}$ against the paper covariance by $\sim1$ unit. The integrated agreement does not mean the bins agree: the LightGBM seed-1 and exact-split central values differ by a median of \SI{0.97}{\percent} per bin, $1.3$ statistical standard deviations (seed~42: \SI{0.98}{\percent}, 1.32; \S\ref{sec:method-config}). Iteration-count effects are negligible,
\begin{center}
\begin{tabular}{@{}lcc@{}}
\toprule
OmniFold configuration & $\chi^{2}/\mathrm{ndf}$, paper covariance & pull SD \\
\midrule
exact-split, 5 iter (cross-check; earlier central value) & $3.66$ & $0.598$ \\
HistGBT, 5 iter (10-seed mean) & $2.70\pm0.04$ & $0.56$ (RMS) \\
LightGBM, 5 iter, seed 42, refined subtraction (quoted) & $2.270$ & $0.535$ \\
LightGBM, 5 iter, seed 42 & $2.654$ & $0.564$ \\
LightGBM, 5 iter, seed 1 & $2.65$ & $0.563$ \\
LightGBM, 10 iter, seed 1 & $2.66$ & $0.565$ \\
\bottomrule
\end{tabular}
\end{center}
All rows except the quoted one use the purity-based subtraction. (i) \textbf{Iteration count, for this statistic}: LightGBM $5\to10$ iters moves the $\chi^{2}$ by $+0.01$ ($2.645\to2.657$). That is the movement of this comparison statistic between five and ten iterations. It does not establish convergence of every functional. (ii) \textbf{The GBDT estimator carries a $\sim1$-unit band}: the exact-split cross-check ($3.66$) is the outlier, while both histogram-binned backends agree at $2.65$--$2.70$. HistGBT, the scikit-learn histogram estimator, is a $10$-seed mean, range $2.63$--$2.74$. The $\pm0.04$ ML-seed jitter is $\sim25\times$ smaller than the $0.96$ exact$\to$hist gap, so these runs differ by backend. That jitter does not measure the exact-split seed spread, which is unmeasured (\S\ref{sec:method}). The exact-gradient trees fit finer shape structure (less implicit smoothing) and so deviate more from the paper's IBU-regularized result. The comparison shows that the $\chi^{2}$ moves with the estimator. It does not separate implicit regularization from a physical data--MC disagreement.
% absence-ok: scoping statement: exact-split seed spread is unmeasured in sec_method.tex
% RE-2 integration W2 (approved, DECISION-20261024-joseph-approve-integration-wording-s1-compute-preserve.md): the seed-42 purity row is E_U (values_table.tsv cfg-lgbm42-chi2 / cfg-lgbm42-pullrms
% purity_value 2.654, 0.564); the quoted row is E_R (refined 2.270, 0.535).
% (\texttt{tension\_iter/})

\paragraph{Conclusion.} For the exact-split cross-check, the tension is a broadly distributed, correlated shape difference concentrated on the low-$\pT$ peak dominated by flux/Muon\_Energy bands; for the LightGBM central value the distance is smaller and a larger share lies in poorly determined directions (\S\ref{sec:results-tension}). Normalization and small-$\lambda$ inversion artifacts do not drive it. A $\sim1$-$\chi^{2}$-unit GBDT estimator/iteration regularization band adds a methodological component. Both point to the same structure: the Bashyal flux\,$\leftrightarrow$\,Muon\_Energy joint block (\S\ref{sec:rank}) is the covariance structure that governs the region carrying the tension.



\subsection{Construction of the two-dimensional covariance}\label{recsec:twod-covariance}

\paragraph{Development-history detail.}
% Moved verbatim from the note's development-history appendix in the compression (DECISION-20261119c-joseph-note-compression-outline.md
% item 3; the director's App. E ruling of 2026-11-19: class H, stopped tests with no result and superseded comparisons; move-manifest-je JE-10).
% [records reorg] former heading: \subsection{Development-history detail moved from the note}

\paragraph{The earlier two-dimensional budget: comparison.}
% [records reorg] former heading: \subsubsection{The earlier two-dimensional budget: comparison}

\paragraph{Comparison with the earlier 2D budget.} The saved products and code of the earlier 2D budget~\cite{TwoDFollowup2026} show the following defects in its construction; the quoted budget corrects each. First, its systematic sweep has a simulated truth input that differs from that of its seed-42 reference, while its saved reconstructed inputs agree. A common displacement cancels to first order in the MAT mean-centered covariance: the mismatch does not invalidate the entire covariance and supplies no correction to subtract, and because the earlier input and code revision are not available, it cannot be attributed more specifically. The sweep used by the quoted universe block does not have this mismatch. Second, the bands Rvn1pi and Rvp1pi encode the same shift under the configured nonresonant-pion tune, and the earlier budget counted it twice; the quoted universe block drops Rvp1pi and counts it once (\S\ref{sec:uq-universes}). The band is small, \SI{0.088}{\percent}, but the double count is a defect in the construction of the covariance; no conclusion about the published analysis's implementation follows. Third, the earlier budget did not propagate background-template statistics, whose error exceeds one tenth of the data's Poisson error in 60 of 205 saved reconstructed bins (a reconstructed-space diagnostic, not a measured unfolded uncertainty); the quoted statistical block draws the background simulation as its own bootstrap stream (\S\ref{sec:uq-boot}). The values printed for the earlier budget were not changed by these findings.
% [records reorg] former heading: \paragraph{Comparison with the earlier 2D budget.}
% (one paragraph since the history restructure's style pass, 2026-11-15: history-cuts/sec_systematics.md S6)
%

\paragraph{Development history (App.~E): detail.}
% Moved verbatim from app_devhistory.tex @ cb76b90d8 in the compression under DECISION-20261120a (qualifier and procedure detail
% may move to the records with a pointer); move-manifest-je JE-17.
% [records reorg] former heading: \subsection{Development history (App.~E): detail moved from the note}

\paragraph{The completeness-divided two-dimensional statistical band.} The undivided band is wider than the toy scatter (pull RMS 0.84, mean pull consistent with zero). The per-bin pull RMS falls as that component's share of the band's variance rises, which corresponds to a data-statistics part about $1.6\times$ wider in $\sigma$ than the toys produce. The toys would scatter too little because their pseudo-data are built on the training MC events. The diagnostic is neither an independent-population coverage test nor a calibration on real data.
% [records reorg] former heading: \subsubsection{The completeness-divided two-dimensional statistical band}

\paragraph{Construction detail of the 2D covariance blocks.}
% Part 3 chunk 2 (the director's Part 3 rulings, 2026-11-19, (a): construction detail to the records under rule 5),
% director-helper-3: moved verbatim from sec_systematics.tex @ lane e09c53477 (draft ae7a78594); rows H3-P3-10..H3-P3-16.
% [records reorg] former heading: \subsection{Construction detail of the 2D covariance blocks (from Part 3, the uncertainty of the 2D cross section)}

\paragraph{PPFX flux universes (from \S\ref{sec:uq-universes}).} The 100 PPFX hadron-production flux universes use the same mean-centered $1/N$ formula.
% [records reorg] former heading: \paragraph{PPFX flux universes (from \S\ref{sec:uq-universes}).}

\paragraph{Flux normalization detail (from \S\ref{sec:uq-universes}, ``Flux $1/\Phi$ normalization'').} The two MINOS muon-energy-scale universes carry the correlated-flux weight, so each is divided by its own flux. That band's integral $\sigma$ is \SI{0.30}{\percent}. The cross term is rank-2; two independent estimators give $\rho=0.87$--$0.90$. Adding this block leaves the systematic rank unchanged because its directions already lie in the universe-covariance span (App.~\ref{sec:rank}). In all 12 playlists the production flux weights equal those of the constraint file and of the earlier 2D budget, with which \S\ref{sec:uq-2d-scope} compares the quoted construction.
% [records reorg] former heading: \paragraph{Flux normalization detail (from \S\ref{sec:uq-universes}, ``Flux $1/\Phi$ normalization'').}
% after "Target-nucleon normalization contributes a fully-correlated rank-1 band of 1.4 %."
% after the sentence on the joint flux and muon-energy-scale block of Bashyal et al.
% the paragraph's last sentence

\paragraph{Bootstrap construction detail (from \S\ref{sec:uq-boot}).} With 300 replicas the Gaussian formula $1/\sqrt{2(N-1)}$ gives a relative standard error of about \SI{4.1}{\percent} on each per-bin $\sigma$, and the per-bin $\sigma$ from a nested half (150) of the replicas is a median 1.000 times that from all 300.
% [records reorg] former heading: \paragraph{Bootstrap construction detail (from \S\ref{sec:uq-boot}).}
% after "... the separable background-only term is contained in the statistical block."

\paragraph{Training-seed block detail (from \S\ref{sec:uq-ml}).} Each setting is a triplet $(s,s{+}1,s{+}2)$ of seeds for the step-1 classifier, step-2 classifier and miss regressor, with $s\in\{42,45,48,51,54\}$; the block has rank~4 (\S\ref{sec:seedscan}). Its per-bin $\sigma$ is a median 0.27 of that of the earlier ten-seed block. The reason is that the bootstrap\,$\times$\,seed interaction, which dominated the earlier scan, is already inside the fixed-seed statistical block. The training-seed block carries it only at its $1/20$ share, a median 0.763 of its diagonal.
% [records reorg] former heading: \paragraph{Training-seed block detail (from \S\ref{sec:uq-ml}).}
% after "A spread of single unfolds on fixed data is a different statistic."
% after "With four degrees of freedom, its variance has a Gaussian relative standard error of about 71 %."

% Part 3 compression at depth (the director, 2026-11-20; DECISION-20261120a item 5), director-helper-3: sentences moved verbatim from the note's Part 3, in order, one paragraph per source paragraph (move manifest H3-P3C-*); revised after cov2d's fact-check and the director's D-item ruling. Where the note keeps text before the moved sentences, the paragraph opens with an italic quote of the paragraph's note sentences that precede them (Grok, codex2). References into the note resolve through main_note.aux (xr). [records reorg] former heading: \section{Part 3 detail moved from the note} [records reorg, not printed] This section holds the construction detail, qualifier detail and secondary checks that the compressed Part~3 no longer prints, as the note printed them. Where the note keeps a sentence of the same paragraph, the paragraph here opens, in italics, with the note section and the note sentences of that paragraph that precede the moved text (or, where none does, the one that follows it); pointers inside those quotes read as in the note.

\paragraph{From the uncertainty of the two-dimensional cross section (\S\ref{sec:systematics}).} The universe, statistical and training-seed blocks are computed with the LightGBM estimator of the quoted 2D central value (\S\ref{sec:method-config}); the universe and statistical blocks use its training seed (seed~42), and the training-seed block uses five seed settings that include it. App.~\ref{app:statmethods} gives the full mathematics and worked examples using the reported numbers (N-D declarations in App.~\ref{app:statmethods-nd}). \emph{After, in \S\ref{sec:uq-universes}, the note's ``Detector and model systematics use the July universe sweep: the 44 error bands and 187 universes of the MINERvA Analysis Toolkit (MAT), identical to the standard systematics set of the published analysis.''} Each universe is re-unfolded with LightGBM at seed~42 and compared with the quoted 2D central value, unfolded with the same estimator and seed. The five lateral bands, which shift reconstructed kinematics rather than event weights, are taken from re-unfolds of re-selected universes (\S\ref{sec:uq-higherdim}) at the same seed. App.~\ref{app:pair} gives the per-band covariance.
% [records reorg] former heading: \subsection{From the uncertainty of the two-dimensional cross section (\S\ref{sec:systematics})}
% From sec_systematics.tex, paragraph 1
% [lead-in, Option B] \emph{After, in \S\ref{sec:systematics}, the note's ``The uncertainty budget sums four blocks: systematic universes, statistical bootstraps, ML training-seed variation, and a run-period block (\S\ref{sec:uq-combine}).'' ``Their independence is an assumption: distinct random streams and physics sources do not by themselves establish it after unfolding.''}
% From sec_systematics.tex, paragraph 2

\paragraph{Flux $1/\Phi$ normalization.} Varying $\Phi$ per flux universe in the dominant $\sigma\propto1/\Phi$ term gives a Flux band of \SI{\fluxBand}{\percent} with the constraint weights, and a combined budget, from our covariance alone, of \SI{\uqMedian}{\percent}. Target-nucleon normalization contributes a fully-correlated rank-1 band of \SI{1.4}{\percent}. \emph{After, in \S\ref{sec:uq-boot}, the note's ``The statistical block jointly bootstraps data and simulation with 300 per-event Poisson(1) weight replicas, each fully re-unfolded with the LightGBM estimator at the central value's training seed (seed~42), and is their sample covariance about the replica mean (the earlier note's bootstrap figure).''} Each replica draws its weights in three independent streams, for data, signal simulation and background simulation; possible shared events between the two simulations were not assessed. Because the background simulation is resampled in every replica, the separable background-only term is contained in the statistical block. The reconstructed-level and truth-level weights of a simulated event share one draw, which preserves within-event migration correlations. Fixing the GBDT seed per replica separates bootstrap and training-seed variances. The quoted central value is one unfold at one seed. Its statistical block is the bootstrap at that seed and already contains the training noise at that seed, so the added component is this covariance of the per-seed means. A spread of single unfolds on fixed data is a different statistic. For two bands whose universes barely change the unfolding inputs, EtaNCEL and NormNCRES, the responses are consistent with instability of the GBDT training contributing at the scale of the training-seed variation (\S\ref{sec:uq-ml}). They neither establish the mechanism for every small band nor justify subtracting an overlap from the budget. The 1N--1P group shows a smaller step, $+\fvval{rp-pc-reco}\,\%$ reconstructed and $+(\fvval{rp-pc-unf}\pm\fvval{rp-pc-unf-se})\,\%$ unfolded, and lies above 1A--1F in every $\pPar$ bin below \SI{8}{GeV/c}, in central value, by \SIrange{\fvval{rp-pc-lev-reco-lo}}{\fvval{rp-pc-lev-reco-hi}}{\percent} at reconstruction level and \SIrange{\fvval{rp-pc-lev-unf-lo}}{\fvval{rp-pc-lev-unf-hi}}{\percent} after unfolding, an excess of unknown cause. Diluted into the combined result, the raw shift of those playlists is small in each bin: every marginal $z$ is below 0.7. The shift's $\pPar$ $\chi^{2}$ against the earlier 2D budget is nevertheless 34.4, where noise would give 2.0. The block is built from separate seed-42 unfolds of three run periods, the playlist groups 1A--1F, 1G--1M and 1N--1P, each unfolded through its own simulated detector response. The block uses each period's unfolded result over that period's own simulated truth. Each period is compared with the exposure-weighted mean of the three. This avoids choosing a reference period. The spread, applied to the central value, gives the block. It is added after the other blocks. Its per-bin median fractional standard deviation is \SI{2.09}{\percent} (84th percentile \SI{3.34}{\percent}, maximum \SI{6.48}{\percent}). The fractional standard deviation is \SIrange{1.17}{3.36}{\percent} on the $\pPar$ projections and \SI{1.97}{\percent} on the integrated cross section. The earlier note's table of the run-period block's effect gives its effect on the total. Taking 1A--1F as the reference would give a smaller block (per-bin median \SI{1.38}{\percent}, integral \SI{1.09}{\percent}). For the same reason, the double-ratio $\chi^{2}$ against any total that contains the block is bounded whatever the data. It does not test the block. The raw 1G--1M $\pPar$ $\chi^{2}$ is 33.8 against the other blocks alone and 2.60 with the block added. These values indicate sensitivity only. 
% G2R: RE-2-shapes run 2 (C_new), g2r-integration/ledger.tsv g2r-cpl-raw-vs-s5, g2r-cpl-raw-vs-q; a source comment citing RE-2 part C for this value gives the earlier value
% [records reorg] former heading: \paragraph{Flux $1/\Phi$ normalization.}
% From sec_systematics.tex, paragraph 4
% [lead-in, Option B] \emph{Before, in \S\ref{sec:uq-universes}, the note's ``The unweighted 100-universe flux covariance has participation-ratio effective rank $\approx1$.''}
% From sec_systematics.tex, paragraph 5
% absence-ok: "not assessed" is the block's own construction scope, unchanged by RE-2; the sentence edited here is W4 (copied with the moved sentence, H3-P3C)
% From sec_systematics.tex, paragraph 7
% [lead-in, Option B] \emph{After, in \S\ref{sec:uq-ml}, the note's ``The training-seed block is the sample covariance, over five training-seed settings, of the mean of 20 bootstrap replicas at each setting.''}
% From sec_systematics.tex, paragraph 9; the paragraph left the note under DECISION-20261120b item 4 (D11, referee).
% From sec_systematics.tex, paragraph 12
% [lead-in, Option B] \emph{After, in \S\ref{sec:uq-combine}, the note's ``The run-period block $C_{\mathrm{PL}}$ is an uncertainty assigned for an observed inconsistency between run periods whose cause is unknown.'' ``The block does not measure a systematic uncertainty from an identified source.'' ``Data taken in different run periods are not mutually consistent within the other uncertainties.'' ``After unfolding, the double ratio (data over simulated truth) of the data-taking periods (playlists) 1G--1M to the earliest periods, 1A--1F, has a $\pPar$ step of $(+6.89\pm0.46)\%$ at \SI{8}{GeV/c}.'' ``The step is already present in the reconstructed data, at $+\fvval{rp-pb-reco}\,\%$ relative to 1A--1F (data minus background over each period's signal simulation, scaled to its exposure; Fig.~\ref{fig:runperiod-steps}(a)).'' ``It therefore survives unfolding at about its reconstruction-level size.''}
% From sec_systematics.tex, paragraph 13
% [lead-in, Option B] \emph{Before, in \S\ref{sec:uq-combine}, the note's ``The block has rank two, and with three period groups its own size is poorly determined.''}
% From sec_systematics.tex, paragraph 14
% [lead-in, Option B] \emph{After, in \S\ref{sec:uq-combine}, the note's ``The block scales the spread of the three periods' unfolded data-to-simulated-truth ratios about their exposure-weighted mean to the central value, so it contains the observed differences by construction (App.~\ref{app:devhistory}), and that containment establishes neither coverage nor whether its size is right; the block's other limits are stated in the records document's two-dimensional limitations.''}

\paragraph{Generator-shape model dependence.} For each of three generators $k$ (GENIE with meson-exchange currents, NuWro and GiBUU), the prior's truth distribution is reweighted by the ratio of that generator's $(\pT,\pPar)$ shape to GENIE's, with its normalization kept. The ratio is constant on each reported $(\pT,\pPar)$ cell. A relative-error gate, fixed before the reweighted unfolds were run, decides each cell's ratio. A cell whose own ratio passes the gate keeps it. Otherwise the cell takes the ratio pooled over a coarser cell, fixed before the measurement, when that pooled error passes the same gate. When neither passes, the cell leaves the prior unchanged. A row outside the grid, an unreported cell, and GiBUU where that sample is not defined also leave the prior unchanged. The data are unfolded again with the central value's estimator; $s_k$ is the per-cell relative change of the result. That sum is a fractional covariance. In cross-section units it is $\mathrm{diag}(x)\,V_{\mathrm{prior}}\,\mathrm{diag}(x)$, with $x$ the cell-integrated cross sections of this result. \emph{Before, in \S\ref{sec:uq-2d-scope}, the note's ``Its median fractional standard deviation over the 205 reported cells is \SI{\cmodelMedian}{\percent} (84th percentile \SI{7.10}{\percent}, maximum \SI{21.10}{\percent}).''} Each generator enters as one fully correlated component, unscaled, so $V_{\mathrm{prior}}$ has rank three.
% [records reorg] former heading: \paragraph{Generator-shape model dependence.}
% From sec_systematics.tex, paragraph 16
% [lead-in, Option B] \emph{After, in \S\ref{sec:uq-2d-scope}, the note's ``The dependence of the result on the generator shape of the prior is reported beside the covariance, as a separate covariance $C_{\mathrm{model}}=\sum_k s_k s_k^{\top}$.''}
% From sec_systematics.tex, paragraph 17

\paragraph{Conventions.} \emph{After, in \S\ref{sec:uq-2d-scope}, the note's ``Following the MAT treatment of universes, each systematic band's covariance is computed about the mean of its universes (constraint-weighted for the flux), with $1/N$ normalization, so it does not include that mean's offset from the central value.''} This follows the MAT treatment of universes (\S\ref{sec:uq-universes}), which we use as a stated convention. The statistical block is the covariance about the replica mean (\S\ref{sec:uq-boot}). Across bands, the per-bin median offset of a band's center from the central value is \SIrange{0.12}{0.57}{\percent} of the central value, or 0.09--2.85 times the band's $\sigma$. 
% G2R: RE-2-shapes run 2 (C_new), g2r-integration/ledger.tsv g2r-centring-min-rel, g2r-centring-max-rel; a source comment citing RE-2 part C for this value gives the earlier value
% [records reorg] former heading: \paragraph{Conventions.}
% From sec_systematics.tex, paragraph 19

\paragraph{Missed events in the universe unfolds.} An earlier construction, since replaced, is described in App.~\ref{app:devhistory-2dchecks}. 
% D1 follow-on (the director's go of 2026-11-16): DECISION-20261114-joseph-d1-g2r-adopted.md items 2-3; COV2D-G2R-DECLARATION.md sec. 10, VALID row (\"removed. The covariance text states that the systematic-universe unfoldings weight each missed event at its universe's truth weight\"). The earlier approximation's two-universe tests are in the ledger: history-cuts/sec_systematics.md S7.
% [records reorg] former heading: \paragraph{Missed events in the universe unfolds.}
% From sec_systematics.tex, paragraph 20
% [lead-in, Option B] \emph{After, in \S\ref{sec:uq-2d-scope}, the note's ``In the systematic-universe unfolds of the adopted covariance, each simulated signal event that fails reconstruction carries that universe's own truth weight.''}

\paragraph{Comparison with the earlier 2D budget.} The quoted budget corrects three defects in the construction of the earlier 2D budget: a simulated-truth input of its systematic sweep that differs from that of its reference, a nonresonant-pion shift counted twice, and background-template statistics that were not propagated. None of the three is present in the quoted budget. The records document gives each defect with its diagnostics.
% [records reorg] former heading: \paragraph{Comparison with the earlier 2D budget.}
% From sec_systematics.tex, paragraph 21

\subsection{Closure, injected-feature and unbinned diagnostics}\label{recsec:twod-validation}

\paragraph{From the validation of the two-dimensional result (\S\ref{sec:validation}).} In a simulated experiment, \emph{truth} is the known target distribution, the \emph{prior} is the starting simulated truth distribution, and the \emph{response} describes reconstruction and selection conditional on truth. Bias is the mean estimate minus that target over repeated experiments. A single closure residual also contains sampling fluctuations. Repeating a numerical calculation checks neither the coverage and width of the intervals nor the adequacy of the physical model.
% [records reorg] former heading: \subsection{From the validation of the two-dimensional result (\S\ref{sec:validation})}
% From sec_validation.tex, paragraph 1
% [lead-in, Option B] \emph{Before, in \S\ref{sec:validation}, the note's ``A closure test compares the unfolded estimate with its known target under a specified deformation.''}

\paragraph{Closure.} \emph{After, in \S\ref{sec:validation}, the note's ``Three families of closure tests use the simulation of playlist ME1A (\S\ref{sec:experiment}), unfolded with LightGBM and five iterations.'' ``Table~\ref{tab:test-index} gives each test's threshold and result.'' ``In the hidden-variable test (a bias in $\Delta\pT=\pT^{\mathrm{reco}}-\pT^{\mathrm{truth}}$, a resolution variable outside the feature set), the leakage is linear in amplitude, $A=0.05/0.10/0.30\to0.57/1.14/3.43\,\%$ (threshold \SI{3.0}{\percent}): the first two amplitudes pass, and the deliberately large $A=0.30$ arm exceeds the threshold.''} The alternative-model test trains on the central-value simulation and takes its pseudo-data from an alternative systematic universe.
% [records reorg] former heading: \paragraph{Closure.}
% From sec_validation.tex, paragraph 2

\paragraph{Injected-feature recovery.} The test requires the residual left after recovering a deviation from the prior to be far smaller than the deviation itself. A data-versus-prior $\chi^2$ with a diagonal covariance of statistical variances only is larger at truth level than at reco level. This is expected at high resolution: detector smearing lowers the reco-level value, and the diagonal covariance omits the unfolding covariance. The comparisons that use a full covariance are the $\chi^2$ against the published result (\S\ref{sec:results-tension}) and the generator comparisons. The earlier note's appendix on the detector-response-mismatch closure defines a separate test, a closure that rescales only the reconstructed $\Eavail$ of the pseudo-data, together with a more complete leading-nucleon follow-up. No result of that reconstruction-only test is reported here; the same appendix reports noise-free closures with a shifted recoil response on the sixteen cells, which have been run. The five-dimensional estimator, including its background subtraction, is evaluated under the technical closure checks of \S\ref{sec:val-fiveaxis}. Its coverage is unvalidated. Each toy draws pseudo-data event by event as Poisson counts with means equal to the simulated reconstructed weights. It resamples the simulation with the per-event Poisson(1) bootstrap of the statistical band and unfolds with that band's estimator, LightGBM with five iterations at a fixed training seed, the learner of the quoted central value at a different seed (\S\ref{sec:method-config}). The reference $T$ is the unfluctuated simulated truth, identical in every toy. The no-fluctuation configuration reproduces both the production closure output and $T$ exactly. In reported bin $b$ of toy $t$, $z_{tb}=(U_{tb}-T_b)/(r_b T_b)$, where $r_b$ is the production band's relative spread. Coverage $C_k$ is the fraction of all (toy, bin) pairs with $|z|\le k$, with intervals from resampling whole toys. A failure at the interim analysis of 100 toys, with 99.5\% intervals, stops the study. The positive control rescales the band on the same toys. At $\times0.7$, $C_2$ moves to $\covTwoDCTwoLowCtl$, below its window. At $\times1.3$, $C_2$ moves to $\covTwoDCTwoHighCtl$, inside its window, and $C_1$ to $0.785$, above its window. The test flags both directions. \emph{After, in \S\ref{sec:coverage-tunes}, the note's ``A second coverage test asks whether the analysis procedure's 2D covariance covers truths held out of its construction.'' ``A pseudo-experiment is inside a cell when the difference between its unfolded result and the held-out truth lies within that cell's standard deviation under the remade covariance.'' ``It is inside the joint region when its $\chi^{2}$ on that covariance, with the degrees of freedom printed below, is at most the $\chi^{2}$ quantile at the one-standard-deviation Gaussian content.'' ``The null-based threshold rule was fixed before any pseudo-experiment was read.''} This containment is why the test is called a coverage test. The test's false-flag rate under its null's law is at most \lrval{cov2d-ar23-size} and \lrval{cov2d-g18-size} per tune. In the same study, the mean $\chi^{2}$ splits into a part from the pseudo-experiments' mean error $\bar e$, $\bar e^{\top}C^{+}\bar e$, and a part from their spread $S$, $\mathrm{tr}(C^{+}S)$, with $C$ the remade covariance. At \fvval{f1-tune-ar23} and \fvval{f1-tune-g18}, the mean-error part is \fvval{t2a-bias-share-ar23}\,\% and \fvval{t2a-bias-share-g18}\,\% of the mean $\chi^{2}$, or \fvval{t2a-bias-ndf-ar23} and \fvval{t2a-bias-ndf-g18} per degree of freedom. The spread part is \fvval{t2a-spread-ndf-ar23} and \fvval{t2a-spread-ndf-g18} per degree of freedom, where 1 would mean that the covariance describes the spread. In the covariance's eigenbasis, the \fvval{t2a-ntop} directions with the largest eigenvalues carry \fvval{t2a-top-ar23}\,\% and \fvval{t2a-top-g18}\,\% of the mean-error part, and the directions beyond the \fvval{t2a-r1} largest carry \fvval{t2a-beyond1-ar23}\,\% and \fvval{t2a-beyond1-g18}\,\%. The decomposition locates the $\chi^{2}$ and leaves its cause open.
% [records reorg] former heading: \paragraph{Injected-feature recovery.}
% From sec_validation.tex, paragraph 3
% [lead-in, Option B] \emph{After, in \S\ref{sec:validation}, the note's ``The 2D truth-reweighting closure with a Gaussian bump passes the recovery test of the unbinned-unfolding guide~\cite{Canelli:2025guide} (Table~\ref{tab:test-index}).''}
% From sec_validation.tex, paragraph 4
% From sec_validation.tex, paragraph 6
% [lead-in, Option B] \emph{After, in \S\ref{sec:validation}, the note's ``None of these tests changes the simulated detector.'' ``The hidden-variable closure does change the effective response: it reweights reconstructed events by a function of $\Delta\pT$ and leaves the truth weights unchanged (\S\ref{sec:closure-hv}).''}
% From sec_validation.tex, paragraph 10
% From sec_validation.tex, paragraph 11
% [lead-in, Option B] \emph{Before, in \S\ref{sec:coverage}, the note's ``The band counts as nominal if its ratio to the true scatter lies in $[0.9,1.1]$, which gives the windows $C_1\in[\covTwoDWinOneLo,\covTwoDWinOneHi]$ and $C_2\in[\covTwoDWinTwoLo,\covTwoDWinTwoHi]$.''}
% From sec_validation.tex, paragraph 13
% From sec_validation.tex, paragraph 17

\paragraph{Unbinned diagnostic (C2ST).} On this training sample, that drop is a discrimination diagnostic. A valid null would require re-running the full cross-fitted OmniFold pipeline per pseudoexperiment. Calibrated unbinned goodness-of-fit remains open~\cite{Canelli:2025guide}.
% [records reorg] former heading: \paragraph{Unbinned diagnostic (C2ST).}
% From sec_validation.tex, paragraph 20
% [lead-in, Option B] \emph{After, in the earlier note's section on other validation diagnostics, the note's ``In the classifier two-sample test (C2ST), a classifier is trained to separate the data from the (reweighted) simulation in the full unbinned feature space, and its held-out area under the ROC curve (AUC) or accuracy is compared with the chance value $0.5$.'' ``The AUC drops from $\approx0.535$ before unfolding to $\approx0.501$ after.''}
